% ID: SEL-D0015-BC-E05
% ESTADO: validada. S016. Revisión matemática, didáctica y visual.
\textbf{Recuentos disjuntos.}\par\nopagebreak[4]
Sean $A$ y $B$ leer el primero y el segundo libro. Los cuatro grupos son: sólo primero $46-16=30$; sólo segundo $34-16=18$; ambos $16$; ninguno $120-(30+18+16)=56$. Suman $120$.

\textbf{a) Alguno de los dos.}\par\nopagebreak[4]
La unión es inclusiva; por inclusión-exclusión,
\[
\boxed{P(A\cup B)=\frac{46+34-16}{120}=\frac{64}{120}=\frac8{15}.}
\]
\textbf{b) Ninguno.}\par\nopagebreak[4]
\[
\boxed{P(A^c\cap B^c)=1-\frac8{15}=\frac7{15}.}
\]
\textbf{c) Solamente el primero.}\par\nopagebreak[4]
\[
\boxed{P(A\cap B^c)=\frac{30}{120}=\frac14.}
\]
\textbf{d) Primero, sabiendo que no leyó el segundo.}\par\nopagebreak[4]
Hay $120-34=86$ estudiantes que no leyeron el segundo; entre ellos $30$ sí leyeron el primero:
\[
\boxed{P(A\mid B^c)=\frac{30/120}{86/120}=\frac{15}{43}.}
\]
La condicionada usa $86$ como recuento de referencia, no $120$. No se supone independencia entre lecturas.
